\section{Visão Geral}

A campanha \texttt{full\_run} possuía 1116 execuções nominais. O analisador
consolidou 1115 pastas com logs suficientes e classificou 847 execuções como
aprovadas pelos critérios automáticos. A Tabela~\ref{tab:resumo-suites} resume os
resultados agregados mais relevantes.

\begin{table}[htbp]
\centering
\small
\setlength{\tabcolsep}{4pt}
\renewcommand{\arraystretch}{1.12}
\begin{tabular}{p{0.25\linewidth}|p{0.18\linewidth}|p{0.17\linewidth}|p{0.3\linewidth}}
\hline
Suíte e transporte & Aprovadas & Entrega & Leitura principal \\
\hline
\texttt{baseline}, broadcast & 150/150 & 100,00\% & Referência mais estável, sem perdas na métrica principal. \\
\hline
\texttt{multiple}, broadcast & 127/150 & 99,76\% & Suporta quatro IDs por sensor, mas expõe perdas pequenas. \\
\hline
\texttt{multiple}, multicast & 119/150 & 91,24\% & Promissor para múltiplos receptores, ainda frágil em descoberta e memória. \\
\hline
\texttt{real\_time}, broadcast & 56/149 & 91,77\% & Sem batching, filas e retries passam a limitar o sistema. \\
\hline
\texttt{active\_filter}, broadcast & 75/90 & 99,99\% & A filtragem funciona, mas a entrega par a par ainda varia. \\
\hline
\texttt{mixed\_frequency}, broadcast & 15/18 & 99,99\% & Frequências diferentes coexistem, com falhas ligadas sobretudo à geração local. \\
\hline
\end{tabular}
\caption{Resumo agregado das suítes executadas na campanha \texttt{full\_run}.}
\label{tab:resumo-suites}
\end{table}

O resultado central é que a variante broadcast com batching foi a única
configuração plenamente estável na matriz de referência. Com
\texttt{linger\_ms = 100}, a suíte \texttt{baseline} manteve 100\% de entrega em
todas as topologias e em frequências de 10 Hz a 1000 Hz. Foram analisados mais de
2,6 milhões de contadores sem perdas na métrica principal. A latência de
aplicação, medida por \texttt{LAT\_RTT}, ficou tipicamente em poucos
milissegundos, com mediana próxima de 3 ms entre 10 Hz e 200 Hz e próxima de
5 ms em 750 Hz e 1000 Hz.

\section{Efeito do Batching}

A suíte \texttt{real\_time} mostrou que o bom desempenho da linha de base depende
diretamente do agrupamento de amostras. Com \texttt{linger\_ms = 0}, cada amostra
é enviada em um quadro ESP-NOW independente. A entrega agregada do broadcast caiu
para 91,77\%, com 57970 quedas na fila do produtor e 231 falhas por limite de
retry. O sistema ainda se manteve quase estável em 10 Hz, mas perdeu robustez a
partir de 200 Hz e não teve execuções aprovadas em 1000 Hz.

Esse comportamento confirma que o WCAN não deve ser tratado como substituto sem
fio de uma rede CAN de controle em tempo real. Sua aplicação adequada é a
aquisição de dados com frequências conhecidas, na qual uma pequena janela de
agrupamento reduz o custo fixo de rádio, callbacks, filas e confirmações.

\section{Múltiplos Fluxos e Filtros}

Na suíte \texttt{multiple}, o broadcast manteve 99,76\% de entrega agregada mesmo
com quatro CAN IDs por sensor e mais de 4,3 milhões de contadores analisados. As
falhas apareceram como pequenas lacunas em receptores específicos ou como poucos
contadores ausentes em todos os receptores esperados. A primeira situação é
compatível com a semântica de primeiro ACK do broadcast, enquanto a segunda indica
perda efetiva, descarte em fila ou limite de retry.

A suíte \texttt{active\_filter} confirmou que receptores configurados com
allowlist não entregam nem confirmam mensagens fora de interesse. A entrega
agregada foi 99,99\%, mas apenas 75 de 90 execuções passaram por todos os
critérios. Em topologias com múltiplos receptores, um ACK válido prova que algum
receptor interessado recebeu o pacote, mas não que todos os receptores de interesse
receberam todos os contadores. Essa é uma limitação esperada do modelo adotado.

\section{Comportamento do Multicast}

O multicast foi avaliado principalmente na suíte \texttt{multiple}. A entrega
agregada de 91,24\% mostra que a estratégia é viável como direção de evolução,
mas ainda não está madura. Em casos representativos, sensores permaneceram por
muito tempo em modo de descoberta, registrando ausência de assinantes vivos mesmo
após receptores observarem pacotes iniciais. Sob carga, esse estado aumentou a
pressão sobre filas e heap, gerando erros como
\texttt{ESP\_ERR\_ESPNOW\_NO\_MEM} e, em alguns casos, reinicialização por falha
de alocação.

A investigação dos logs não apontou incompatibilidade de canal como causa
principal. O firmware fixa o canal ESP-NOW em 1 e não foram observados erros
característicos de canal incorreto. As falhas são melhor explicadas por
convergência incompleta de assinaturas, pressão de memória e falta de uma política
mais forte para períodos sem assinantes válidos.

\section{Limitações dos Ensaios}

Os resultados foram obtidos em bancada, com placas alimentadas e monitoradas por
USB. Essa configuração é adequada para comparar variantes e encontrar gargalos de
firmware, mas ainda não valida vibração, alimentação real, posicionamento de
antenas, interferência eletromagnética do carro, massa final do encapsulamento ou
injeção contínua no barramento CAN do veículo.

Também foram observados sinais de risco em memória e pilha. O menor high-water
mark de pilha observado nas agregações principais do broadcast foi 64 bytes, e os
logs multicast mostraram pressão de heap em períodos prolongados de descoberta.
Esses pontos não invalidam a linha de base, mas devem ser tratados antes de uma
validação embarcada prolongada.

\section{Rastreabilidade dos Requisitos}

A Tabela~\ref{tab:rastreabilidade-requisitos} resume o estado dos requisitos ao
final dos ensaios. O resultado mais forte é a validação em bancada do broadcast
com batching. Os requisitos ligados a múltiplos receptores, integração embarcada
e robustez física permanecem parciais.

\begin{table}[htbp]
\centering
\scriptsize
\setlength{\tabcolsep}{3pt}
\renewcommand{\arraystretch}{1.12}
\begin{tabular}{p{0.22\linewidth}|p{0.54\linewidth}|p{0.16\linewidth}}
\hline
Requisito & Evidência principal & Status \\
\hline
Desempenho e latência & Broadcast com batching entregou 100\% na suíte \texttt{baseline}, de 10 Hz a 1000 Hz. & Atendido em bancada \\
\hline
Conexão instantânea & Broadcast transmite sem associação prévia; multicast ainda falha quando a descoberta não converge. & Parcial \\
\hline
Múltiplas conexões & Topologias até cinco placas foram testadas; primeiro ACK não garante entrega a todos os receptores. & Parcial \\
\hline
Transparência CAN & O CAN ID é preservado e a fronteira de integração com TWAI foi definida, mas os testes finais usaram log serial. & Parcial \\
\hline
Requisitos físicos e custo & ESP32 reduz componentes e custo, mas placa final, fixação e massa ainda não foram validadas. & Parcial \\
\hline
Modularidade & Papéis por configuração serial, filtros por CAN ID e múltiplas frequências foram exercitados nos testes. & Atendido em bancada \\
\hline
\end{tabular}
\caption{Estado dos requisitos após a campanha de testes.}
\label{tab:rastreabilidade-requisitos}
\end{table}

\section{Conclusão dos Resultados}

Com os dados da campanha \texttt{full\_run}, a configuração mais defensável para
uso controlado é o broadcast com batching habilitado. Ela atende parcialmente ao
objetivo do trabalho ao permitir aquisição sem fio em bancada com semântica
próxima à do CAN, na qual uma mensagem é considerada confirmada quando algum
receptor interessado a reconhece. Para cenários em que todos os receptores devem
receber todos os pacotes, será necessário amadurecer o multicast ou introduzir
confirmação explícita por receptor.
